\par\Needspace{9\baselineskip}\bigskip
\pdfbookmark[3]{Exposition, [22]}{pallas-II-article22}
\begin{center}[22.]\end{center}
\begingroup
\setlength{\lineskip}{4pt}\setlength{\lineskiplimit}{4pt}
\noindent\hspace*{\parindent}Il est évident que les variations instantanées de chaque élément seront les produits de \(\mathrm{d}t\) en des fonctions du tems et des élémens de la planète \SourceLeaf{472}perturbée et de la planète perturbante. Dans la première approximation, tous ces élémens sont supposés constans et le tems n’y entre qu’en tant qu’il se trouve dans les valeurs des longitudes moyennes des deux planètes, lesquelles seront
\[
\begin{aligned}
L&=N+nt\\
L'&=N'+n't,
\end{aligned}
\]
en distinguant par des accens toutes les quantités qui se rapportent à la planète perturbante. De cette manière la variation instantanée de chaque élément aura la forme \(\mathrm{d}t.(\text{fonction de }L\text{ et }L')\). Il est de plus évident que cette fonction sera \textit{périodique} relativement à \(L\text{ et }L'\) parcequ’elle ne changera [pas] de valeur en variant \(L\text{ et }L'\) de \(360^{0}\) ou d’un multiple de \(360^{0}\). Mais on sait que des fonctions de ce genre peuvent être exprimées par une suite composée de termes de la forme \(k\cos(iL+i'L')\) ou \(k\sin(iL+i'L')\), où \(k\) est un coëfficient constant, \(i\) désigne tous les nombres entiers positifs (zéro compris), \(i'\) tous les nombres entiers positifs ou négatifs. Le terme, où \(i=i'=0\), donnera dans l’intégration \(kt\) ce qui est la variation séculaire de l’élément; les autres termes donneront des parties périodiques de la forme \(\dfrac{k}{in+i'n'}\sin(iL+i'L')\). Il faudrait excepter les termes où \(in+i'n'=0\), s’il s’en trouvait, lesquels ajouteraient \(kt\cos(iL+i'L')\) à la variation séculaire, mais il est évident qu’il ne peut y en avoir à moins que les mouvemens moyens n’aient une raison rationelle. Enfin il faut ajouter une constante.
\par\endgroup
